% ------------------------------------------------------------------------------
In addition to the basic \tool{sbatch} options presented earlier, there are a 
few additional options that are generally useful:

\begin{itemize}
\item
%\texttt{-m bea}: requests that the scheduler e-mail you when a job (b)egins;
%(e)nds; (a)borts. Mail is sent to the default address of,
%\texttt{"username@encs.concordia.ca"}, unless a different address is supplied (see, 
%\texttt{-M}). The report sent when a job ends includes job 
%runtime, as well as the maximum memory value hit (\api{maxvmem}). 
\texttt{--mail-type=TYPE}: requests that the scheduler e-mail you when a job changes
state. Where \texttt{TYPE} is \texttt{ALL}, \texttt{BEGIN}, \texttt{END}, or \texttt{FAIL}.
% TODO: verify
Mail is sent to the default address of, \\
\texttt{"<ENCSusername>@encs.concordia.ca"}, which you can consult
via \texttt{webmail.encs} via the VPN, on login.encs via \tool{alpine}
or setup forwarding to @concordia.ca address or offsite,
unless a different address is supplied (see, \texttt{--mail-user}).
% TODO: double-check
The report sent when a job ends includes job 
runtime, as well as the maximum memory value hit (\api{maxvmem}). 

\item
%\texttt{-M email@domain.com}: requests that the scheduler use this e-mail 
%notification address, rather than the default (see, \texttt{-m}). 
\texttt{--mail-user email@domain.com}: requests that the scheduler use this e-mail 
notification address, rather than the default (see, \texttt{--mail-type}). 

\item
%\texttt{-v variable[=value]}: exports an environment variable that can be used by the script.
\texttt{--export=[ALL | NONE | variables]}: exports environment variable(s) that can be used by the script.

\item
%\texttt{-l h\_rt=[hour]:[min]:[sec]}: sets a job runtime of HH:MM:SS. Note 
%that if you give a single number, that represents \emph{seconds}, not hours. 
\texttt{-t [min]} or \texttt{DAYS-HH:MM:SS}: sets a job runtime of min or HH:MM:SS. Note 
that if you give a single number, that represents \emph{minutes}, not hours. 

\item
%\texttt{-hold\_jid [job-ID]}: run this job only when job [job-ID] finishes. Held jobs appear in the queue. 
\texttt{--depend=[state:job-ID]}: run this job only when job [job-ID] finishes. Held jobs appear in the queue. 

\end{itemize}

The many \tool{sbatch} options available are read with, \texttt{man sbatch}. Also 
note that \tool{sbatch} options can be specified during the job-submission 
command, and these \emph{override} existing script options (if present). The 
syntax is, \texttt{sbatch [options] PATHTOSCRIPT}, but unlike in the script, 
the options are specified without the leading \verb+#SBATCH+
(e.g., \texttt{sbatch -J sub-test --chdir=./ --mem=1G ./tcsh.sh}).